\documentclass[letterpaper, 11pt, leqno]{article}
\usepackage{paper}
\usepackage{math}
\usepackage{exam}
\hypersetup{pdftitle={Problem Set on Differential Equations}}

\begin{document}

\title{Problem Set on Differential Equations}
\author{Pascal Michaillat}
\date{}
\paperurl{https://pascalmichaillat.org/x/}
\begin{titlepage}
\maketitle
\end{titlepage}

\section*{Problem 1}

Find the solution of the initial value problem 
\begin{align*}
\dot{a}(t) &=r a(t) +s \\
a(0) &=a_{0}
\end{align*}
where $r$ and $s$ are known constants.

\section*{Problem 2}

Find the solution of the initial value problem 
\begin{align*}
\dot{a}(t) &=r(t) a(t) +s(t) \\
a(0) &=a_{0}
\end{align*}
where $r(t)$ and $s(t)$ are known functions of time $t$.

\section*{Problem 3}

Consider the linear system of differential equations given by
\begin{equation*}
\dot{\bm{x}}(t)=\bs{
\begin{array}{cc}
1 & 1 \\ 
4 & 1
\end{array}} \bm{x}(t).
\end{equation*}
\begin{enumerate}
\item Find the general solution of the system.
\item What would you need to find a specific solution of the system?
\item Draw the trajectories of the system.
\end{enumerate}

\section*{Problem 4}

Consider the initial value problem 
\begin{align*}
\dot{k}(t) &=s f(k(t)) -\d k(t) \\
k(0) &=k_{0}
\end{align*}
where the saving rate $s\in (0,1)$, the capital depreciation rate $\d \in (0,1)$, and the production function $f$ satisfies the \textit{Inada conditions}. That is, $f$ is continuously differentiable and 
\begin{align*}
f(0)&=0\\
f'(x)&>0\\
f''(x)&<0\\
\lim_{x\to 0} f'(x)&=+\infty\\
\lim_{x\to +\infty} f'(x)&=0.
\end{align*}

\begin{enumerate}
\item Give a production function $f$ that satisfies the Inada conditions.
\item Find the steady state of the system.
\item Draw the dynamic path of $k(t)$ and show that it converges to the steady state.
\end{enumerate}

\section*{Problem 5}

The solution of the problem studied in Problem 4 is characterized by a system of two nonlinear first-order differential equations:
\begin{align*}
\dot{k}(t) &=f(k(t)) -c(t)-\d k(t)\\
\frac{\dot{c}(t)}{c(t)} &=\a  a k(t)^{\a -1}-(\d +\r).
\end{align*}
The first differential equation is the law of motion of capital. The second differential equation is the Euler equation, which describes the optimal path of consumption over time.

\begin{enumerate}
\item Draw the phase diagram of the system.
\item Linearize the system around its steady state.
\item Show that the steady state is a saddle point locally.
\item Suppose the economy is in steady state at time $t_{0}$ and there is an unanticipated decrease in the discount factor $\r$. Show on your phase diagram the transition dynamics of the model.
\end{enumerate}

\section*{Problem 6}

Consider the following system of two nonlinear first-order differential equations associated with an investment problem:
\begin{align*}
\dot{k}(t) &=\bp{\frac{q(t)-1}{\c}} k(t) \\
\dot{q}(t) &=r q(t)-f'(k(t)) -\frac{1}{2 \c} (q(t)-1)^{2}.
\end{align*}
The first differential equation is the law of motion of capital $k(t)$. The second differential equation is the law of motion of the costate variable $q(t)$.

\begin{enumerate}
\item Draw the phase diagram.
\item Show that the steady state is a saddle point locally.
\end{enumerate}

\section*{Problem 7}

Consider a discrete-time version of the typical growth model:
\begin{align*}
k(t+1) &=f(k(t)) -c(t) + (1-\d) k(t) \\
c(t+1) &=\b \bs{1+f'(k(t)) -\d} c(t) .
\end{align*} 
The discount factor $\b \in (0,1)$, the rate of depreciation of capital $\d \in (0,1)$, initial capital $k_{0}$ is given, and the production function $f$ satisfies the Inada conditions. These two equations are a system of first-order difference equations. Whereas a system of first-order differential equations relates $\dot{\bm{x}}(t) $ to $\bm{x}(t)$, a system of first-order difference equations relates $\bm{x}(t+1)$ to $\bm{x}(t)$.

We will see that we can study a system of first-order difference equations with the tools that we used to study systems of first-order differential equations. In particular, we can use phase diagrams to understand the dynamics of the system.

\begin{enumerate}
\item Construct a phase diagram for the system. First, define 
\begin{align*}
\D k & \equiv k(t+1) -k(t), \\
\D c & \equiv c(t+1) -c(t).
\end{align*}
Second, draw the $\D k=0$ locus and the $\D c=0$ locus on the $(k,c)$ plane. Finally, find the steady state as the intersection of the $\D k=0$ locus and the $ \D c=0$ locus.
\item Show that the steady state is a saddle point in the phase diagram.
\end{enumerate}

\section*{Problem 8}

We consider the following optimal growth problem. Given initial human capital $h_{0}$ and initial physical capital $k_{0}$, choose consumption $c(t)$ and labor $l(t) $ to maximize utility
\begin{equation*}
\int_{0}^{\infty}e^{-\r t} \ln{c(t)} dt
\end{equation*}
subject to
\begin{align*}
\dot{k}(t) &=y(t)-c(t)-\d k(t) \\
\dot{h}(t) &=b [1-l(t)] h(t).
\end{align*}
Output $y(t)$ is given by
\begin{equation*}
y(t) = a k(t)^{\a} [l(t) h(t)]^{\b}.
\end{equation*}
We also impose that $0 \le l(t)\le 1$. The discount factor $\r>0$, the rate of depreciation of physical capital $\d>0$, the constants $a>0$ and $b>0$, and the production function parameters $\a\in (0,1)$ and $\b \in(0,1)$.

\begin{enumerate}
\item Give state and control variables.
\item Write down the present-value Hamiltonian for this problem.
\item Derive the optimality conditions. 
\item Show that the growth rate of consumption $c(t)$ is
\begin{equation*}
\frac{\dot{c}}{c}=\frac{\a y}{k}-(\d +\r).
\end{equation*}
\item From now on, we assume that $b=0$. Show that $l=1$ in optimum.
\item Draw the phase diagram in the $(k,c)$ plane.
\item Show on the phase diagram that the steady state of the system is a saddle point.
\item Derive the Jacobian of the system.
\item Use the Jacobian to show that the steady state of the system is a saddle point.
\end{enumerate}

\end{document}